\section{Analytic Obstructions}
\label{sec:analytic}

\begin{theorem}[Prime-supported Euler products have a natural boundary]
  \label{thm:euler-product-natural-boundary}
  Let \((b_n)_{n\ge 1}\subseteq\ZZ_{\ge 0}\) and assume that
  \(b_p>0\) for infinitely many primes \(p\). Suppose that
  \[
    Z_b(z):=\prod_{n\ge 1}(1-z^n)^{-b_n}
  \]
  converges to a holomorphic function on \(\abs{z}<1\). Then the unit
  circle is a natural boundary for \(Z_b\).
\end{theorem}

\begin{proof}
  If \(b_p>0\), then the factor \((1-z^p)^{-b_p}\) has poles at all
  \(p\)th roots of unity. Since every other factor has non-negative
  exponent, no cancellation can occur there. Thus every prime \(p\)
  with \(b_p>0\) contributes poles at the \(p\)th roots of unity. These
  roots are dense in the unit circle because the set of such primes is
  infinite. A holomorphic continuation across any boundary arc would
  therefore meet infinitely many poles, which is impossible. Hence
  \(\abs{z}=1\) is a natural boundary.
\end{proof}

\begin{proposition}[Imaginary-periodicity rigidity of finite determinants]
  \label{prop:finite-zeta-imaginary-periodicity}
  For every finite matrix \(M\), the function
  \[
    F_M(s):=\det(I-e^{-s}M)^{-1}
  \]
  is \(2\pi i\)-periodic. In particular, the Riemann
  \(\zeta\)-function cannot coincide with \(F_M\) on any open subset of
  \(\Re(s)>1\).
\end{proposition}

\begin{proof}
  The identity \(e^{-(s+2\pi i)}=e^{-s}\) implies
  \(F_M(s+2\pi i)=F_M(s)\). On the other hand, for large real
  \(\sigma>1\),
  \[
    \zeta(\sigma+2\pi i)-\zeta(\sigma)
    =\sum_{n\ge 2}n^{-\sigma}(e^{-2\pi i\log n}-1).
  \]
  The \(n=2\) term has size comparable to \(2^{-\sigma}\), while the
  tail is \(O(3^{-\sigma})\), so the difference is non-zero for
  sufficiently large \(\sigma\). Hence \(\zeta(s)\) is not
  \(2\pi i\)-periodic on \(\Re(s)>1\).
\end{proof}

\begin{remark}
  \Cref{thm:euler-product-natural-boundary,prop:finite-zeta-imaginary-periodicity}
  are complementary. The first obstructs finite-state models at the
  level of the boundary geometry of the Euler product. The second
  obstructs them already after the logarithmic change of variables.
\end{remark}
